\section{Explicit Algebraic Model Duality}
\label{sec:ch20-model-duality}

\subsection{Why the order-parameter exchange needs a map of algebras}
\label{subsec:ch20c-motivation}

In Section~\ref{sec:ch20-formal-vertices}, we defined the charge-density-wave (CDW) and superconducting (SC) vertices as formal series within an uncompressed oscillator--charge algebra. Their evaluated oscillator correlators depend on the reciprocal interaction parameters $g$ and $g^{-1}$, suggesting a duality between the density-wave and pairing channels. However, the similarity between two resulting formulas does not constitute a proof of duality. To establish duality, we construct a single, well-defined algebraic map that consistently transforms all oscillator generators, charge labels, fermionic Klein signs, and ordered position-dependent phases. This explicit map allows us to transport Hamiltonians and states between the two models.

The charge sectors reveal a structural hurdle before analyzing the oscillators. The CDW Klein word shifts the source charge by $e_R-e_L=(1,-1)$. In contrast, our SC word is an annihilation pair, shifting the charge by $-e_R-e_L=(-1,-1)$. A simple reflection of the left charge maps $(1,-1)$ to $(1,1)$, which is the charge shift of the adjoint operator $W_{\mathrm{SC}}^*$, not $W_{\mathrm{SC}}$ itself. Consequently, any algebraic exchange must incorporate an adjoint operation. We will also see that resolving the remaining source-phase mismatch requires a factor of $\zeta^x$. This complexity shows why simply swapping $g$ and $g^{-1}$ is insufficient.

We consider an even finite lattice length $L=2h$ ($h\geq1$), an integer position $x$ modulo $L$, the canonical phase $\zeta=e^{2\pi i/L}$, and a finite oscillator-mode cutoff $1\leq m\leq M$ ($M\geq1$). The left and right chiralities remain ordered as $R<L$. We briefly recall the definitions required to make this duality explicit, relying on the positive polynomial vacuum and formal vertices from Section~\ref{sec:ch20-formal-vertices}.

An \emph{algebra automorphism} is a bijective linear map preserving
products and the unit.  A \emph{star automorphism} also satisfies
$D(a^*)=D(a)^*$.  In particular it does not reverse products:
$D(ab)=D(a)D(b)$.  We shall construct a map of this kind rather
than infer one from its proposed action on a few expectation values.

\subsection{The parameter and oscillator involutions}
\label{subsec:ch20c-oscillators}

Write the real interaction parameters as
\begin{equation}
 P=(u,c,s,g,g^{-1}),\qquad
 c^2-s^2=1,\qquad g=(c-s)^2>0,\qquad
 g^{-1}=(c+s)^2,
 \label{eq:ch20c-params}
\end{equation}
with $u>0$.  In this notation $g^{-1}$ is the reciprocal of $g$,
not an independent coefficient.  If the Hamiltonian is parametrized
by $v_1,v_2$, we use
\begin{equation}
 v_1=u(c^2+s^2),\qquad v_2=2ucs.
 \label{eq:ch20c-velocities}
\end{equation}
The symbol $g$ should not be confused with the Umklapp coefficient
$g_U$ of Chapter~21 or with an energy-budget bound.

\begin{definition}[Interaction-family involution]
\label{def:ch20c-param}
Define
\begin{equation}
 P^*=(u,c,-s,g^{-1},g).
 \label{eq:ch20c-Pstar}
\end{equation}
The star on $P^*$ designates a \emph{parameter involution}; it is
not complex conjugation or an adjoint of a Hamiltonian.
\end{definition}

\begin{lemma}[Parameter exchange]
\label{lem:ch20c-param}
The map $P\mapsto P^*$ preserves $u$ and $v_1$, reverses $v_2$,
exchanges $g$ with $g^{-1}$, and satisfies $(P^*)^*=P$.
\end{lemma}
\begin{proof}
Since $c^2-(-s)^2=1$, the transformed record is admissible.
Furthermore
\[
 (c-(-s))^2=(c+s)^2=g^{-1},\qquad
 (c+(-s))^2=(c-s)^2=g.
\]
In~\eqref{eq:ch20c-velocities}, $s$ enters $v_1$ only through
$s^2$ and enters $v_2$ linearly.  Changing its sign therefore gives
$v_1^*=v_1$ and $v_2^*=-v_2$.  Changing the sign twice returns the
initial record.  The reciprocal identity also follows directly from
$(c-s)(c+s)=c^2-s^2=1$.
\end{proof}

Let $\mathfrak A_{\mathrm{osc}}$ be the unital complex star-algebra
generated by $A_{\nu m},C_{\nu m}$ ($\nu\in\{R,L\}$,
$1\leq m\leq M$), with
\begin{align}
 [A_{\nu m},C_{\eta n}]&=m\delta_{\nu\eta}\delta_{mn}I,
 & [A_{\nu m},A_{\eta n}]&=0,
 & [C_{\nu m},C_{\eta n}]&=0,\nonumber\\
 A_{\nu m}^*&=C_{\nu m},& C_{\nu m}^*&=A_{\nu m}.
 \label{eq:ch20c-ccr}
\end{align}
It is a quotient of a free word algebra by the two-sided star-stable
ideal generated by these relations, not a finite-dimensional matrix
algebra.  Its familiar polynomial representation has $C_{\nu m}$
acting by multiplication by $X_{\nu m}$ and $A_{\nu m}$ by
$m\partial/\partial X_{\nu m}$.

\begin{proposition}[Oscillator sign automorphism]
\label{prop:ch20c-Dosc}
There is a unique complex-linear star automorphism
$D_{\mathrm{osc}}$ of $\mathfrak A_{\mathrm{osc}}$ specified by
\begin{align}
 D_{\mathrm{osc}}(A_{Rm})&=-A_{Rm},&
 D_{\mathrm{osc}}(C_{Rm})&=-C_{Rm},\nonumber\\
 D_{\mathrm{osc}}(A_{Lm})&= A_{Lm},&
 D_{\mathrm{osc}}(C_{Lm})&= C_{Lm}.
 \label{eq:ch20c-Dosc}
\end{align}
It is involutive: $D_{\mathrm{osc}}^2=\mathrm{id}$.
\end{proposition}
\begin{proof}
Define the substitution on arbitrary finite words multiplicatively
and on finite linear combinations linearly.  For the $R$ generators,
\[
 [-A_{Rm},-C_{Rn}]=[A_{Rm},C_{Rn}]
 =m\delta_{mn}I,
\]
and for $L$ nothing changes.  Mixed-chirality commutators are zero
before and after the substitution; the vanishing $AA$ and $CC$
commutators are likewise preserved.  Because both members of each
adjoint pair acquire the same real sign,
$D_{\mathrm{osc}}(A_{\nu m}^*)
 =D_{\mathrm{osc}}(A_{\nu m})^*$.
Thus every defining relation maps into the relation ideal and the
substitution descends to the quotient.  Repeating it changes each
$R$ generator by two minus signs and fixes every $L$ generator,
so the induced map is its own inverse.
\end{proof}

In the notation of Section~\ref{sec:ch20-formal-vertices}, define the
Bogoliubov automorphism $\beta_P$ by
\begin{align}
 \widetilde C_{Rm}(P)&=cC_{Rm}+sA_{Lm},&
 \widetilde A_{Rm}(P)&=cA_{Rm}+sC_{Lm},\nonumber\\
 \widetilde C_{Lm}(P)&=cC_{Lm}+sA_{Rm},&
 \widetilde A_{Lm}(P)&=cA_{Lm}+sC_{Rm}.
 \label{eq:ch20c-dressed}
\end{align}
All these equalities are inside the uncompressed algebra.  They
must not be read as identities of compressed finite-CAR matrices.

\begin{lemma}[Intertwining with Bogoliubov dressing]
\label{lem:ch20c-intertwine}
On the entire oscillator star algebra,
\begin{equation}
 D_{\mathrm{osc}}\circ\beta_P
 =\beta_{P^*}\circ D_{\mathrm{osc}}.
 \label{eq:ch20c-intertwine}
\end{equation}
In particular,
\begin{align}
 D_{\mathrm{osc}}\widetilde C_{Rm}(P)
   &=-\widetilde C_{Rm}(P^*),&
 D_{\mathrm{osc}}\widetilde A_{Rm}(P)
   &=-\widetilde A_{Rm}(P^*),\nonumber\\
 D_{\mathrm{osc}}\widetilde C_{Lm}(P)
   &=\widetilde C_{Lm}(P^*),&
 D_{\mathrm{osc}}\widetilde A_{Lm}(P)
   &=\widetilde A_{Lm}(P^*).
 \label{eq:ch20c-dressed-transform}
\end{align}
\end{lemma}
\begin{proof}
For the right creation generator,
\[
 D_{\mathrm{osc}}(cC_{Rm}+sA_{Lm})
   =-cC_{Rm}+sA_{Lm}
   =-(cC_{Rm}-sA_{Lm}),
\]
which is $\beta_{P^*}(-C_{Rm})$.  For the left creation generator,
\[
 D_{\mathrm{osc}}(cC_{Lm}+sA_{Rm})
   =cC_{Lm}-sA_{Rm}=\beta_{P^*}(C_{Lm}).
\]
The annihilation formulas follow by taking adjoints.  Since the two
sides of~\eqref{eq:ch20c-intertwine} are algebra homomorphisms,
agreement on generators proves agreement on every finite word and
hence on the quotient.
\end{proof}

Define the Hermitian dressed chiral fields with their opposite
Fourier orientations,
\begin{align}
 Q_R(P,x)&=\sum_{m=1}^{M}\frac{
  \zeta^{-mx}\widetilde C_{Rm}(P)
 +\zeta^{mx}\widetilde A_{Rm}(P)}{m},\nonumber\\
 Q_L(P,x)&=\sum_{m=1}^{M}\frac{
  \zeta^{mx}\widetilde C_{Lm}(P)
 +\zeta^{-mx}\widetilde A_{Lm}(P)}{m},
 \label{eq:ch20c-Q}\end{align}
\begin{equation}
 X_{\mathrm{CDW}}(P,x)=(c-s)(Q_R(P,x)+Q_L(P,x)),\quad
 X_{\mathrm{SC}}(P,x)=(c+s)(Q_R(P,x)-Q_L(P,x)).
 \label{eq:ch20c-X}
\end{equation}
These are well-defined finite sums of oscillator words, even though
their formal exponentials have infinitely many coefficients.
\begin{corollary}[Exchange of oscillator phases]
\label{cor:ch20c-X}
For all integer positions $x$,
\begin{align}
 D_{\mathrm{osc}}X_{\mathrm{CDW}}(P,x)
  &=-X_{\mathrm{SC}}(P^*,x),\nonumber\\
 D_{\mathrm{osc}}X_{\mathrm{SC}}(P,x)
  &=-X_{\mathrm{CDW}}(P^*,x).
 \label{eq:ch20c-Xtransform}
\end{align}
\end{corollary}
\begin{proof}
By~\eqref{eq:ch20c-dressed-transform},
$D_{\mathrm{osc}}Q_R(P,x)=-Q_R(P^*,x)$ and
$D_{\mathrm{osc}}Q_L(P,x)=Q_L(P^*,x)$.
For example,
\[
 D_{\mathrm{osc}}X_{\mathrm{CDW}}(P,x)
 =(c-s)(-Q_R(P^*,x)+Q_L(P^*,x)).
\]
The SC coefficient at $P^*$ is $c+(-s)=c-s$, giving the first
identity.  The second uses $c-(-s)=c+s$.  No sign is supplied by
complex conjugation: these are linear algebra substitutions.
\end{proof}

\subsection{Reflecting charge and preserving the Klein signs}
\label{subsec:ch20c-charge}

The charge representation occupies the algebraic direct sum
$\mathcal C=\mathbb C^{(\mathbb Z^2)}$ of finite-support vectors.
Its distinguished orthonormal kets are $\delta_N$ with
$N=(N_R,N_L)$.  Fix the same integer half-band offset $h$ as in
Section~\ref{sec:ch20-formal-vertices}.  Write $\mathsf N_\nu$ for the diagonal charge operator and
$N_\nu$ for the integer coordinate of a source ket.  The charge
operators and root-of-unity phases act by
\begin{equation}
 \mathsf N_\nu\delta_N=N_\nu\delta_N,
 \qquad Z_\nu(x)\delta_N=\zeta^{N_\nu x}\delta_N.
 \label{eq:ch20c-NZ}
\end{equation}
The signed lowering operators and their adjoints are
\begin{align}
 F_\nu\delta_N&=P_\nu(N)\delta_{N-e_\nu},\nonumber\\
 F_\nu^*\delta_N&=P_\nu(N+e_\nu)\delta_{N+e_\nu},\nonumber\\
 P_R(N)&=(-1)^{h+N_R-1},\qquad
 P_L(N)=(-1)^{N_R+N_L-1}.
 \label{eq:ch20c-Klein}
\end{align}
The left sign follows from the ordered-phase convention:
\[
 P_L(N)=(-1)^{(h+N_R)+(h+N_L-1)}.
\]
The integer $2h$ drops out of parity.  Thus $F_L$ contains a sign
depending on the \emph{right} charge; this convention cannot be
suppressed without changing the product coefficients.

\begin{definition}[Charge reflection and induced map]
\label{def:ch20c-charge}
Set
\begin{equation}
 f(N_R,N_L)=(N_R,-N_L),\qquad
 U_f\delta_N=\delta_{f(N)}.
 \label{eq:ch20c-f}
\end{equation}
The permutation $U_f$ is an involutive isometry of the algebraic
inner-product space $\mathcal C$.  On the \emph{generated adjointable}
star algebra $\mathfrak A_{\mathrm{ch}}$ of the operators in
\eqref{eq:ch20c-NZ}--\eqref{eq:ch20c-Klein}, define
\begin{equation}
 D_{\mathrm{ch}}(B)=U_fBU_f^{-1}=U_fBU_f.
 \label{eq:ch20c-Dcharge}
\end{equation}
The map is not defined here as a star automorphism of the entire
space of arbitrary endomorphisms of $\mathcal C$.
\end{definition}

\begin{lemma}[Exact action on the charge generators]
\label{lem:ch20c-charge}
The map $D_{\mathrm{ch}}$ preserves the generated star algebra and
satisfies
\begin{align}
 D_{\mathrm{ch}}(\mathsf N_R)&=\mathsf N_R,&
 D_{\mathrm{ch}}(\mathsf N_L)&=-\mathsf N_L,\nonumber\\
 D_{\mathrm{ch}}(Z_R(x))&=Z_R(x),&
 D_{\mathrm{ch}}(Z_L(x))&=Z_L(x)^*,\nonumber\\
 D_{\mathrm{ch}}(F_R)&=F_R,&
 D_{\mathrm{ch}}(F_L)&=-F_L^*.
 \label{eq:ch20c-charge-images}
\end{align}
It is a complex-linear, involutive star automorphism.
\end{lemma}
\begin{proof}
For each diagonal operator the claim follows by applying
$U_fBU_f$ to $\delta_N$: its eigenvalue is evaluated at $f(N)$.
In particular $\zeta^{-N_Lx}=\overline{\zeta^{N_Lx}}$, proving
the left-phase identity.  For the right signed shift,
\[
 U_fF_RU_f\delta_N=P_R(fN)\delta_{N-e_R}
 =P_R(N)\delta_{N-e_R}=F_R\delta_N.
\]
For the left shift, its reflected action is
\[
 U_fF_LU_f\delta_N
  =P_L(fN)\delta_{N+e_L},
\]
while the adjoint acts with coefficient $P_L(N+e_L)$.  Their ratio is
\[
 \frac{P_L(fN)}{P_L(N+e_L)}
 =\frac{(-1)^{N_R-N_L-1}}{(-1)^{N_R+N_L}}
 =(-1)^{-2N_L-1}=-1.
\]
Thus $D_{\mathrm{ch}}(F_L)=-F_L^*$, including the minus sign.
Every displayed image is again in the generated star algebra,
so conjugation restricts to an endomorphism of it.  As $U_f^2=I$,
this endomorphism is invertible and involutive.  Finally,
$U_f$ preserves the algebraic inner product, hence
$(U_fBU_f)^*=U_fB^*U_f$ on finite-support kets.  This proves
star preservation without assuming that the whole endomorphism
algebra carries an appropriate adjoint operation.
\end{proof}

\subsection{Ordered zero-mode words and exact vertex exchange}
\label{subsec:ch20c-vertices}

The CDW and pairing Klein--phase factors, in precisely the ordering
used for the finite-CAR coefficients and the formal vertices, are
\begin{align}
 K_{\mathrm{CDW}}(x)
 &=F_R^*F_L\,\zeta^{(\mathsf N_L-\mathsf N_R-1)x},
 \label{eq:ch20c-Kcdw}\\
 K_{\mathrm{SC}}(x)
 &=F_RF_L\,\zeta^{(\mathsf N_R+\mathsf N_L)x}.
 \label{eq:ch20c-Ksc}
\end{align}
The rightmost character acts on the \emph{source} ket before the
first shift.  This is why the power of $\zeta$ is not obtained by
simply replacing $N_L$ with $-N_L$ in a final answer.

\begin{proposition}[Klein--phase transformation with source shifts]
\label{prop:ch20c-Ktransform}
With $R<L$, the canonical root $\zeta$, and the signs
\eqref{eq:ch20c-Klein}, one has the exact operator identities
\begin{align}
 D_{\mathrm{ch}}K_{\mathrm{CDW}}(x)
 &=\zeta^x K_{\mathrm{SC}}(x)^*,\nonumber\\
 D_{\mathrm{ch}}K_{\mathrm{SC}}(x)
 &=\zeta^x K_{\mathrm{CDW}}(x)^*.
 \label{eq:ch20c-Kdual}
\end{align}
The scalar $\zeta^x$ and both adjoints are indispensable.
\end{proposition}
\begin{proof}
First record the action of the original words on any source
$N=(N_R,N_L)$:
\begin{align}
 K_{\mathrm{CDW}}(x)\delta_N
 &=a_C(N)\,\zeta^{(N_L-N_R-1)x}
       \delta_{N+e_R-e_L},\nonumber\\
 a_C(N)&=P_L(N)P_R(N-e_L+e_R),
 \label{eq:ch20c-actionC}\\
 K_{\mathrm{SC}}(x)\delta_N
 &=a_S(N)\,\zeta^{(N_R+N_L)x}
       \delta_{N-e_R-e_L},\nonumber\\
 a_S(N)&=P_L(N)P_R(N-e_L).
 \label{eq:ch20c-actionS}
\end{align}
The coefficients $a_C,a_S$ are real signs.  Therefore their
adjoint actions are obtained by reversing the charge shifts and
conjugating the \emph{phase evaluated at the original source}.

Apply the reflected CDW word to $\delta_N$.  Its amplitude is
$a_C(fN)\zeta^{(-N_L-N_R-1)x}$, and its target is
$N+e_R+e_L$.  On the other hand, $K_{\mathrm{SC}}(x)^*$ has
this same target, because its inverse source is
$N+e_R+e_L$.  The corresponding signed coefficient is
$a_S(N+e_R+e_L)$ and its character is
$\zeta^{-(N_R+N_L+2)x}$.  Direct use of
\eqref{eq:ch20c-Klein} gives
\begin{align*}
 a_C(fN)
 &=(-1)^{N_R-N_L-1}(-1)^{h+N_R}
 =(-1)^{h+N_L+1},\\
 a_S(N+e_R+e_L)
 &=(-1)^{N_R+N_L+1}(-1)^{h+N_R}
 =(-1)^{h+N_L+1}.
\end{align*}
The signs match, and their characters differ by precisely
$\zeta^x$:
\[
 \zeta^{(-N_L-N_R-1)x}
 =\zeta^x\zeta^{-(N_R+N_L+2)x}.
\]
Hence $D_{\mathrm{ch}}K_{\mathrm{CDW}}(x)
=\zeta^xK_{\mathrm{SC}}(x)^*$ on every basis vector.

For the reflected SC word, the target is $N-e_R+e_L$.
The inverse source for the CDW adjoint is then $N-e_R+e_L$.
The signs are
\begin{align*}
 a_S(fN)&=(-1)^{N_R-N_L-1}(-1)^{h+N_R-1}
 =(-1)^{h+N_L},\\
 a_C(N-e_R+e_L)
 &=(-1)^{N_R+N_L-1}(-1)^{h+N_R-1}
 =(-1)^{h+N_L},
\end{align*}
while the phases obey
\[
 \zeta^{(N_R-N_L)x}
 =\zeta^x\zeta^{-(N_L-N_R+1)x}.
\]
This proves the second operator identity on every $\delta_N$.
Because these vectors span $\mathcal C$, basis extensionality
finishes the proof.  Reordering the Klein operators or evaluating
the diagonal phase after a shift would spoil this calculation.
\end{proof}

Set $\mathfrak A=\mathfrak A_{\mathrm{ch}}\otimes_{\mathbb C}
\mathfrak A_{\mathrm{osc}}$.  Since charge and oscillator generators
act on different tensor factors, they commute in $\mathfrak A$.
Define
\begin{equation}
 D=D_{\mathrm{ch}}\otimes D_{\mathrm{osc}}:
  \mathfrak A\longrightarrow\mathfrak A.
 \label{eq:ch20c-Dtensor}
\end{equation}
On pure tensors $D(B\otimes a)=D_{\mathrm{ch}}(B)\otimes
D_{\mathrm{osc}}(a)$; bilinearity ensures that this is
well-defined on the algebraic tensor product.  It is a unital,
involutive star automorphism by the two preceding constructions.

For the indeterminate $t$ assume $t$ is central and $t^*=t$.
Extend $D$ coefficientwise to the noncommutative formal-series algebra
$\mathfrak A[[t]]$ by
\[
 D\left(\sum_{n\geq0}a_nt^n\right)
  =\sum_{n\geq0}D(a_n)t^n.
\]
This preserves multiplication: the coefficient of $t^n$ in a
product is the \emph{finite} sum
$\sum_{j=0}^{n}a_jb_{n-j}$; applying $D$ distributes over the
sum and preserves the order of each product.  It preserves star,
and its square is the identity.  No convergence or substitution
$t=1$ is used.

Recall that the Weyl-type vertices are, with the formal exponential
$\operatorname{Exp}_{\mathrm f}(tY)
=\sum_{n\geq0}Y^nt^n/n!$,
\begin{equation}
 W_J(P,t,x)=\frac1L K_J(x)
   \operatorname{Exp}_{\mathrm f}(itX_J(P,x)),
 \qquad J\in\{\mathrm{CDW},\mathrm{SC}\}.
 \label{eq:ch20c-W}
\end{equation}
These are not normal-ordered operators and their formal parameter
is not an analytic coupling.

\begin{theorem}[Exact coefficientwise CDW--pairing duality]
\label{thm:ch20c-Wdual}
For every admissible $P$, finite $M$, integer $x$, and formal $t$,
\begin{align}
 D W_{\mathrm{CDW}}(P,t,x)
 &=\zeta^xW_{\mathrm{SC}}(P^*,t,x)^*,\nonumber\\
 D W_{\mathrm{SC}}(P,t,x)
 &=\zeta^xW_{\mathrm{CDW}}(P^*,t,x)^*.
 \label{eq:ch20c-vertex-duality}
\end{align}
Both equalities hold in $\mathfrak A[[t]]$ coefficient by
coefficient.  In particular they assert no value at $t=1$.
\end{theorem}
\begin{proof}
For a Hermitian phase $X$ and $t^*=t$,
\[
 \operatorname{Exp}_{\mathrm f}(itX)^*
 =\operatorname{Exp}_{\mathrm f}(-itX)
\]
coefficient by coefficient: taking adjoints of $X^n$ reverses the
word order but leaves it equal to $X^n$, and conjugates $i^n$.
By Corollary~\ref{cor:ch20c-X},
\[
 D\operatorname{Exp}_{\mathrm f}
       (itX_{\mathrm{CDW}}(P,x))
 =\operatorname{Exp}_{\mathrm f}
       (-itX_{\mathrm{SC}}(P^*,x)).
\]
By Proposition~\ref{prop:ch20c-Ktransform}, the full CDW image is
\[
 \frac{\zeta^x}{L} K_{\mathrm{SC}}(x)^*
 \operatorname{Exp}_{\mathrm f}
       (-itX_{\mathrm{SC}}(P^*,x)).
\]
The charge word commutes with the oscillator series.  Thus this is
$\zeta^xW_{\mathrm{SC}}(P^*,t,x)^*$, with the \emph{reversed}
order required by taking the adjoint of $W_{\mathrm{SC}}$.
The second equality is identical with CDW and SC exchanged.
\end{proof}

\subsection{Hamiltonian covariance and transported states}
\label{subsec:ch20c-Hstate}

An identity of interaction families becomes stronger when the same
map intertwines their Hamiltonians and prescribed states.  Both
statements require an explicit choice of Hamiltonian and of the
source charge: they are not assertions about every finite lattice
Hamiltonian or every fixed-charge state.

\begin{theorem}[Symmetric Hamiltonian covariance]
\label{thm:ch20c-H}
Assume the symmetric, chemically adjusted zero-mode convention and
real $v_1,v_2$ as in~\eqref{eq:ch20c-velocities}.  In the
uncompressed oscillator--charge algebra define
\begin{align}
 H(P)={}&\frac{2\pi}{L}\sum_{m=1}^{M}\Bigl[
 v_1\bigl(C_{Rm}A_{Rm}+C_{Lm}A_{Lm}\bigr)
 +v_2\bigl(C_{Rm}C_{Lm}+A_{Rm}A_{Lm}\bigr)\Bigr]
 \nonumber\\
 &+\frac{\pi}{L}\Bigl[
 v_1(\mathsf N_R^2+\mathsf N_L^2)
 +2v_2\mathsf N_R\mathsf N_L\Bigr].
 \label{eq:ch20c-H}
\end{align}
Then, with $v_1(P^*)=v_1(P)$ and $v_2(P^*)=-v_2(P)$,
\begin{equation}
 D H(P)=H(P^*).
 \label{eq:ch20c-Hdual}
\end{equation}
\end{theorem}
\begin{proof}
For each $m$, both same-chirality number words are unchanged:
\[
 D(C_{Rm}A_{Rm})=(-C_{Rm})(-A_{Rm})=C_{Rm}A_{Rm},
\]
and the $L$ word is fixed.  Every mixed-chirality pairing word
contains exactly one $R$ factor, hence
\[
 D(C_{Rm}C_{Lm})=-C_{Rm}C_{Lm},\qquad
 D(A_{Rm}A_{Lm})=-A_{Rm}A_{Lm}.
\]
Thus the oscillator quadratic expression retains its form with
$v_2$ replaced by $-v_2$.  For the zero-mode polynomial,
\[
 D(\mathsf N_R^2+\mathsf N_L^2)
 =\mathsf N_R^2+\mathsf N_L^2,\quad
 D(\mathsf N_R\mathsf N_L)=-\mathsf N_R\mathsf N_L.
\]
Applying these identities term by term to the finite sum gives
\eqref{eq:ch20c-Hdual}.  No commutation of a creator past an
annihilator was needed, so the given normal and pairing order was
preserved.
\end{proof}

The corresponding states must be transported along the charge
permutation.  Recall from Section~\ref{sec:ch20-formal-vertices}
that the bare oscillator vacuum $\omega_0$ is represented by the
constant polynomial $1$ and that
$\omega_P=\omega_0\circ\beta_P^{-1}$.  For a source label $N$, the
full state $\omega_{P,N}$ is the vector functional at
$\delta_N\otimes1$ in the tensor product representation with the
dressed oscillator vacuum.  This is a normalized positive
\emph{complex-valued} functional on $\mathfrak A$; its formal-series
extension will be denoted $\widehat\omega_{P,N}$.

\begin{theorem}[Charge-indexed state transport]
\label{thm:ch20c-state}
For every $a\in\mathfrak A$ and every $N\in\mathbb Z^2$,
\begin{equation}
 \omega_{P^*,f(N)}(D a)=\omega_{P,N}(a).
 \label{eq:ch20c-statedual}
\end{equation}
The identity extends coefficientwise to all
$A(t)\in\mathfrak A[[t]]$:
\begin{equation}
 \widehat\omega_{P^*,f(N)}(D A(t))
   =\widehat\omega_{P,N}(A(t)).
 \label{eq:ch20c-formal-statedual}
\end{equation}
It is a \emph{family} identity; it does not claim that an arbitrary
single fixed-charge state is invariant under $D$.
\end{theorem}
\begin{proof}
First show that the polynomial vacuum is invariant under the bare
oscillator sign change.  In its polynomial representation define
$V_R$ by
\[
 (V_Rp)(X_{Rm},X_{Lm})
   =p(-X_{Rm},X_{Lm})\qquad(1\leq m\leq M),
\]
where the notation abbreviates the simultaneous change of sign of
all right-mode variables.  Every monomial is multiplied by $\pm1$;
therefore $V_R$ is an involutive isometry for the positive weighted
monomial Gram form, and $V_R1=1$.  Directly on polynomials,
\[
 V_R\pi_0(a)V_R^{-1}=\pi_0(D_{\mathrm{osc}}a)
\]
for the multiplication and differentiation generators, hence for
every oscillator word.  Consequently
\[
 \omega_0(D_{\mathrm{osc}}a)
 =\langle1,V_R\pi_0(a)V_R1\rangle
 =\omega_0(a).
\]
By Lemma~\ref{lem:ch20c-intertwine} and invertibility of the
Bogoliubov maps,
\[
 \beta_{P^*}^{-1}D_{\mathrm{osc}}
  =D_{\mathrm{osc}}\beta_P^{-1}.
\]
Thus
\begin{align*}
 \omega_{P^*}(D_{\mathrm{osc}}a)
 &=\omega_0\bigl(\beta_{P^*}^{-1}
                     D_{\mathrm{osc}}a\bigr)\\
 &=\omega_0\bigl(D_{\mathrm{osc}}\beta_P^{-1}a\bigr)
 =\omega_P(a).
\end{align*}

On the charge side the permutation itself implements
$D_{\mathrm{ch}}$, so for each charge word $B$,
\begin{align*}
 \langle\delta_{f(N)},D_{\mathrm{ch}}(B)
              \delta_{f(N)}\rangle
 &=\langle U_f\delta_N,U_fBU_f\,U_f\delta_N\rangle\\
 &=\langle\delta_N,B\delta_N\rangle.
\end{align*}
For a pure tensor $B\otimes a$ the full expectation is the
product of these two factor expectations.  Combining the identities
proves~\eqref{eq:ch20c-statedual} on pure tensors and, by
linearity, on finite sums of pure tensors, which exhaust the
algebraic tensor product.  Applying the established equality to
one coefficient at a time yields
\eqref{eq:ch20c-formal-statedual}; this uses no notion of positivity
for a general formal power series.
\end{proof}

The Hamiltonian and state laws together may be summarized as
\[
 H(P)\ \xrightarrow{\ D\ }\ H(P^*),\qquad
 (P,N)\ \xrightarrow{\ D\ }\ (P^*,fN).
\]
Notice that the second arrow refers to a \emph{pair of labels}, not
to an invariant state at $N$.  In particular, for $N_L\neq0$,
$fN\neq N$.  The map on operators is an automorphism of the
uncompressed algebra; its existence alone does not construct a
finite-energy-budget isomorphism of the Fock subspaces from
Section~\ref{sec:ch20-projected-products}.

\begin{remark}[Chemical potentials must be transformed or adjusted]
\label{rem:ch20c-mu}
An unadjusted zero mode can contain linear terms.  For example, if
\[
 H_{\mu}(P)=H(P)-\mu(\mathsf N_R+\mathsf N_L),
\]
then
\[
 D H_{\mu}(P)
 =H(P^*)-\mu(\mathsf N_R-\mathsf N_L)
 =H_{\mu}(P^*)+2\mu\mathsf N_L.
\]
The simple covariance law fails for nonzero $\mu$ on the full
charge algebra.  One may instead use the chemically adjusted
symmetric zero-mode convention above, restrict to a sector where
the residual vanishes, or explicitly reflect the coefficient of the
left linear term in a larger Hamiltonian family.  Such adjustments
must be stated; one cannot silently discard a chemical-potential
contribution.  Additive scalar offsets likewise require consistent
parameter transformation if they are included.
\end{remark}

\subsection{A four-site calculation}
\label{subsec:ch20c-example}

\begin{example}[Duality with signs, phases, and a nonzero coupling]
\label{ex:ch20c-L4}
Take $L=4$, $h=2$, $M=1$, $x=1$, $\zeta=i$, and the hyperbolic
parameters
\[
 u=8,\qquad c=\frac54,\qquad s=\frac34.
\]
Indeed $c^2-s^2=(25-9)/16=1$, while
\[
 (v_1,v_2)=(17,15),\quad
 (g,g^{-1})=\left(\frac14,4\right),\quad
 (v_1^*,v_2^*)=(17,-15).
\]
Thus this example is interacting at the quadratic level: $v_2$
does not vanish.  Let the charge source be $N=(0,0)$.
The ordered Klein signs are
\[
 P_L(0,0)=-1,\qquad P_R(1,-1)=+1,
\]
and the CDW source phase is $\zeta^{-1}=-i$.  Accordingly
\[
 K_{\mathrm{CDW}}(1)\delta_{(0,0)}
 =(-i)(-1)(+1)\delta_{(1,-1)}
 =i\delta_{(1,-1)}.
\]
Applying the charge reflection to its initial and final kets gives
\[
 D_{\mathrm{ch}}K_{\mathrm{CDW}}(1)\delta_{(0,0)}
  =i\delta_{(1,1)}.
\]
To check the right-hand side of the duality law independently,
$K_{\mathrm{SC}}(1)$ sends $\delta_{(1,1)}$ to
$\delta_{(0,0)}$ with coefficient $+1$: its Klein signs multiply
to $-1$ and its source phase is $i^{2}=-1$.  Hence
\[
 iK_{\mathrm{SC}}(1)^*\delta_{(0,0)}
 =i\delta_{(1,1)},
\]
in exact agreement.  Omitting the factor $\zeta^x=i$ would
produce an incorrect complex coefficient, and omitting the adjoint
would produce the opposite charge shift.

The oscillator transformation is equally explicit:
\[
 D_{\mathrm{osc}}\widetilde C_{R1}(P)
 =-\frac54 C_{R1}+\frac34 A_{L1}
 =-\widetilde C_{R1}(P^*),
\]
with the left dressed generator mapped without an overall minus
sign.  Therefore
\(D X_{\mathrm{CDW}}(P,1)=-X_{\mathrm{SC}}(P^*,1)\),
so the entire formal vertex identity, not only its constant
coefficient, follows coefficient by coefficient.

For a distinct illustration of the Hamiltonian statement, choose
$N=(1,2)$.  Its zero-mode energy is
\[
 E_0(P;1,2)
 =\frac\pi4[17(1^2+2^2)+30(1)(2)]
 =\frac{145\pi}{4}.
\]
The reflected charge is $(1,-2)$.  With $v_2^*=-15$,
\[
 E_0(P^*;1,-2)
 =\frac\pi4[17(1^2+(-2)^2)+2(-15)(1)(-2)]
 =\frac{145\pi}{4}.
\]
These are equal because the charge \emph{and} the coupling are
reflected.  They are not a demonstration of spectral equivalence of
finite-band fermionic Hamiltonians.
\end{example}

\subsection{Logical limits of the duality}
\label{subsec:ch20c-limits}

The preceding proofs establish exact algebraic identities at any
finite $M$ on the \emph{uncompressed} oscillator--charge model,
and coefficientwise identities for its formal vertex series.  They
do not establish the following further statements.

\begin{enumerate}
\item \emph{Finite-CAR equivalence.}  The finite fermionic Hilbert
space has Pauli constraints and the energy-budget projections of
Section~\ref{sec:ch20-projected-products}.  An isomorphism between
those budget carriers, compatible with their source, intermediate,
and target sectors, requires its own construction.  The intermediate
projection residuals derived there remain present.  In particular,
$D$ is not being asserted as a unitary conjugation on the finite-CAR
Fock space.

\item \emph{Finite-budget spectral equivalence.}  Even if two
quadratic expressions have the same transformation law in the
uncompressed algebra, compressing them with projections that are not
transported by an established intertwiner can change their
spectra.  A statement of unitary equivalence of compressed
Hamiltonians requires a specified map between budget subspaces and
a proved intertwining equality on those spaces.

\item \emph{Continuum or compactification duality.}  The algebraic
exchange $g\leftrightarrow g^{-1}$ echoes the familiar
$\phi\leftrightarrow\theta$ interchange in Luttinger-liquid
language (see, for context, \cite{giamarchi2004}).  With the bare
field convention $\phi=Q_R+Q_L$ and
$\theta=Q_R-Q_L$, the sign change of the right component produces
$\phi\mapsto-\theta$ and $\theta\mapsto-\phi$ (with the corresponding
dressed-family comparison).  A continuum $T$-duality, however,
needs a constructed continuum theory, compactification radius and
charge lattice, and an appropriate treatment of vertex sectors and
boundary conditions; none follows from the finite-mode star map.

\item \emph{Analytic vertices or physical phases.}  The maps on
$\mathfrak A[[t]]$ are coefficientwise.  Neither setting $t=1$ nor
claiming power-law asymptotics, a phase transition, or a gap is
justified without convergence estimates or a separate limiting and
spectral argument.  Nor does the positive state
$\omega_{P,N}$ automatically identify with a ground-state
expectation of an interacting finite-CAR Hamiltonian.
\end{enumerate}

In summary, the duality involution has been explicitly verified on the formal algebra constructed here. We have mapped the bare generators, their dressed families, the signed charge shifts, the ordered characters, the finite polynomial Hamiltonians, the charge-indexed positive states, and the coefficientwise Weyl-type vertices. Crucially, this algebraic model duality is not a statement of full interacting-CAR equivalence. Extending this formal mapping into a proven unitary symmetry of the finite-lattice CAR Hamiltonian, or evaluating it as a continuum dynamical equivalence, remains an open problem. As throughout Chapter 20, these derivations are unverified mathematical drafts awaiting formalization in Lean.

The models discussed in this chapter incorporate only two species of fermions (right and left movers). Physical models often include spin, which doubles the mode content to four species. In Chapter~21, we will motivate this four-species extension, establishing the species ordering, charge shifts, and parity conventions needed to define spinful interacting systems without ambiguity.
